% =====================================================================
%  EXAMPLE. Base resume, role family: backend, market: intl, language: en
%  Built by the build-resume skill from example/facts.md.
%  Fictional person and companies.
% =====================================================================
\documentclass[a4paper,10pt]{article}

\usepackage[T1]{fontenc}
\usepackage[utf8]{inputenc}
\usepackage[english]{babel}
\usepackage{geometry}
\usepackage{enumitem}
\usepackage{hyperref}
\usepackage{xurl}
\usepackage{microtype}
\usepackage{titlesec}

\geometry{margin=1.35cm}
\setlength{\parindent}{0pt}
\setlist[itemize]{leftmargin=1.15em, itemsep=2pt, topsep=2pt, parsep=0pt}

\titleformat{\section}{\large\bfseries}{}{0em}{}[\titlerule]
\titlespacing{\section}{0pt}{6pt}{4pt}

\hypersetup{colorlinks=true, urlcolor=black, linkcolor=black}
\Urlmuskip=0mu plus 1mu
\emergencystretch=2em
\sloppy
\input{glyphtounicode}
\pdfgentounicode=1

% \mbox prevents a break after "/": otherwise extraction yields "CI/ CD".
\newcommand{\CICD}{\mbox{CI/CD}}

% ---------- DATA (profile.json, CONTACT-INTL set) ----------
\newcommand{\CVName}{Alexey Morozov}
\newcommand{\CVTitle}{Backend Engineer (Python / Go)}
\newcommand{\CVEmail}{alexey.morozov@example.org}
\newcommand{\CVPhone}{+7 900 000 33 44}
\newcommand{\CVTelegram}{amorozovdemo}

\begin{document}

% FACT: IDENT-name CONTACT-INTL TITLE-backend IDENT-telegram LOC-ekb LANG-en-b2
{\LARGE \textbf{\CVName}}\\
\textbf{\CVTitle}\\
Email: \href{mailto:\CVEmail}{\CVEmail} \quad
Phone: \CVPhone \quad
Telegram: \href{https://t.me/\CVTelegram}{@\CVTelegram}\\
Remote or hybrid \quad English --- B2

\section*{Professional Summary}

% FACT: SUM-backend-en
Backend engineer with 5 years of experience, 3 of them in fintech. I build and run
microservices in Python (FastAPI, asyncio) and Go: the payment gateway my team owns
handles around 400 requests per second at peak. Particular strength in PostgreSQL
beyond CRUD --- query plans, indexing and partitioning. I take services through to
production rather than to merge: tests, metrics, and on-call incident analysis.

\section*{Technical Skills}

% FACT: SKL-lang SKL-frameworks SKL-db SKL-infra SKL-arch SKL-obs SKL-test
\textbf{Languages:} Python (5 years, asyncio, typing), Go (2 years, services in
production), SQL, Bash, Git; Rust --- basic\\
\textbf{Frameworks and protocols:} FastAPI, aiohttp, Celery, gRPC, REST API;
Django --- 1.5 years\\
\textbf{Databases and queues:} PostgreSQL (query plans, indexing, partitioning),
Redis (caching, distributed locks); Kafka --- producing and consuming in production
tasks; MongoDB --- basic\\
\textbf{Infrastructure:} Docker, GitLab \CICD{} (\mbox{Continuous Integration/Delivery}),
Linux; Kubernetes --- deploying via existing manifests, debugging pod failures\\
\textbf{Architecture and reliability:} microservice decomposition, idempotency,
retries, partial-failure handling; saga pattern --- limited exposure\\
\textbf{Observability and testing:} Prometheus, Grafana, Sentry, structured logging,
pytest, testcontainers

\section*{Experience}

% FACT: EXP-aurora
\textbf{Aurora Fintech} \hfill \textit{March 2023 --- Present}\\
\textit{Backend Engineer}
\begin{itemize}
    \item Build and maintain a payment gateway in Python (FastAPI, asyncio): payment
          intake, provider reconciliation, retries and operation idempotency.
    \item Designed the idempotency scheme for payment operations, eliminating double
          charges on retried requests.
    \item Ported two internal services from Python to Go, roughly halving response
          time on heavy endpoints.
    \item Tuned PostgreSQL: query plans, missing indexes, monthly partitioning of the
          transactions table.
    \item Set up Prometheus and Grafana metrics and alerts for my services; on-call
          rotation, incident analysis and postmortems.
    \item Integration and unit testing with pytest and testcontainers across critical paths.
    \item Code review and mentoring: two interns, both hired after their internship.
\end{itemize}

% FACT: EXP-cityline
\textbf{Cityline Retail} \hfill \textit{July 2021 --- February 2023}\\
\textit{Backend Engineer}
\begin{itemize}
    \item Developed catalogue and cart services for an e-commerce platform in Python
          (Django, later FastAPI), exposing REST APIs to web and mobile clients.
    \item Moved price-list export to asynchronous tasks (Celery, Redis), unblocking
          the main request path and fitting the nightly window.
    \item Built a Redis caching layer for the catalogue, cutting peak load on the
          primary database.
    \item Extracted the catalogue service during a monolith split, including data migration.
    \item Maintained GitLab \CICD{} pipelines: tests, linters, image builds.
\end{itemize}

% FACT: EXP-cityline-intern
\textbf{Cityline Retail} \hfill \textit{February 2021 --- June 2021}\\
\textit{Software Engineering Intern}
\begin{itemize}
    \item Fixed bugs in an existing codebase, wrote unit tests, learned to navigate
          a large legacy code base. Received a full-time offer after the internship.
\end{itemize}

\section*{Projects}

% FACT: PRJ-ratelimit
\textbf{Rate limiter in Go} \hfill \textit{personal, open source}
\begin{itemize}
    \item Rate-limiting library: token bucket and sliding window, Redis-backed state
          for distributed mode, benchmarks and race-condition tests.
    \item Adopted by one of the internal services at work.
\end{itemize}

% FACT: PRJ-migrator
\textbf{Database schema migration tool} \hfill \textit{internal library at work}
\begin{itemize}
    \item CLI on top of Alembic: ordered migration runs across several services,
          rollback, and pre-deploy checks for incompatible changes.
    \item Extracted into an internal library, used by three teams.
\end{itemize}

\section*{Education}

% FACT: EDU-utu
\textbf{Ural Technical University} --- B.Sc. Software Engineering (2017--2021)

% FACT: EDU-go EDU-pg
\textbf{Additional training}
\begin{itemize}
    \item Go for backend engineers --- ported two internal services from Python to Go
          as the practical outcome.
    \item PostgreSQL query plans and optimisation --- slow query log analysis, query
          rewriting, index selection.
\end{itemize}

\section*{Awards}

% FACT: AWD-hackathon
\begin{itemize}
    \item Fintech Hackathon 2022 --- 2nd place, team category.
\end{itemize}

\end{document}
